\section{  Relationship between Normal Subgroups and Their Images}
Consider a group $G$ that has a normal subgroup $N$ and $H$ and that $H \subseteq N$. 
\begin{theorem}
Because $N\supseteq H$, both $N$ and $H$ are normal subgroups of $G$. Thus $G/N$ is the homomorphic image of $G/H$ under the natural homomorphsim
$\pi: G/H \goes G/N$ defined by 
$(xG)\pi = xN$. 
\end{theorem}
\begin{proof}

1) Define a function $\phi: G\goes G/N$ by
\[ x\phi = xN.\]
This is the standard canonical projection map. It is a surjective group homomorphism with $ker(\phi) = N$. 

2) We want to construct a map $\overline{\phi}: G/H \goes G/N$.
Define $\overline{\phi}$ by:
\[ (xH) \overline{\phi} = x\phi = xN.\]

3) We must show that if $xH = yH$, then $(xH)\overline{\phi} = (yH)\overline{\phi}.$
\begin{eqnarray*}
xH &=& yH \goes y^{-1}x \in H.\cr
H &\subseteq& N \mbox{  given}\cr
y^{-1}x &\in& N\cr
xN &=& N\cr
(xH)\overline{\phi} = (yH)\overline{\phi}\cr
\end{eqnarray*}
The map is well-defined because $ H \subseteq ker(\phi)$. 
\end{proof}

4) $\overline{\phi}$ is a Homomorphism:
\[ (xH) (yH) \overline{\phi} = (xyH)\overline{\phi} = xyN = (xN)(yN) = (xH)\overline{\phi} (yH)\overline{\phi}.\]

5) $\overline{\phi}$ is Surjective:
Let $yN \in G/N$. 
Since $y\in G$, we can form the coset $yH \in G/H$:
\[ (yH)\overline{\phi} = yN.\] Every element in the co-domain has a pre-image
 
6) Apply the First Isomorphism Theorem
By the First Isomorphism Theorem, the domain modulo the kernel is isomorphic to the image

\[ (G/H)/ker(\overline{\phi}) \cong Im(\overline{\phi}) = G/N .\]
Since $G/N$ is isomorphic to a quotient group of $G/H$, it is by definition a homomorphic image of $G/H$.\qed \\

{\bf Lattice Isomorphism Theorems}
 
\begin{theorem} [First Isomorphism Theorem:] Let $\phi$ be a group homomorphism from $G$ to $G'$. Then the mapping from $G/Ker(\phi)$ to $G\phi$ given by 
$g Ker(\phi) \goes g\phi$ is an isomorphism.
\end{theorem}
\begin{example}Consider the homomorphism $\phi: \Z \goes \Z_4$ such that $x\phi = x\mod4$.
In this case let $G$ be the group of integers (under addition) \[ G = \{..., -4, -3, -2, -1, 0, 1, 2, 3, 4, ...\}\] 
and \[ G' = \Z_4 = \{ 0, 1, 2, 3\}. \]
Also \[ Ker(\phi) = \{ ...-8, ,-4, 0, 4, 8, ...\} = \langle4\rangle \mbox{   (the cyclic subgroup generated by 4)} \] 
So $G\phi = \{ 0, 1, 2, 3\}$. Note that $G\phi$ doesn't have to be the entire group $G'$. In this example it is the case that $G\phi = G'$.
Then \[ G/ Ker(\phi) = \{ 0 + \langle4\rangle, 1 + \langle4\rangle, 2 + \langle4\rangle, 3+\langle4\rangle \} \mbox{   (four cosets)}\]
Notice that we can create a mapping $\overline{\phi}$ that mapps $G/Ker(\phi)$ into $G\phi$ as follows: 
\[ \overline{\phi}: G/Ker(\phi) \goes g\phi, \] where
\begin{eqnarray*}
\overline{\phi}(0+\langle4\rangle) &=& (0)\phi = 0\cr
\overline{\phi}(1+\langle4\rangle) &=& (1)\phi = 1\cr
\overline{\phi}(2+\langle4\rangle) &=& (2)\phi = 2\cr
\overline{\phi}(3+\langle4\rangle) &=& (3)\phi = 3
\end{eqnarray*}
\end{example}

\begin{theorem} [Second Isomorphism Theorem:] The Second Isomorphism Theorem: If  $S$
 is a subgroup of $G$  and $N$  is a normal subgroup of $G$ , then the product $SN$ is a subgroup, and the quotient 
$SN / N$  is isomorphic to  $S / (S \cap N$.\end{theorem}

\begin{theorem}[Third Isomorphism Theorem:] Normal Subgroups of Quotients: If $N$ and $H$ are normal subgroups of $G$ (where $N \subseteq H$), then 
$H/N$ is  is a normal subgroup of $G/N$. Furthermore, the second quotient is 
isomorphic to the primary quotient: \[ (G / N )/  (H / N)  \cong G / H. \]
\end{theorem}
The Third Isomorphism Theorem is analogous to fractions:
\[ {G / N \over H / N} \cong {G\over H}.\]

\begin{theorem} [Fourth Isomorphism Theorem:]
There is a direct, structure-preserving correspondence between the subgroups of a quotient group $G / N$
 and the subgroups of $G$
 that contain $N$.
\end{theorem}

\begin{theorem}
The mapping $\phi: \Z\goes\Z_n$ defined by $m\phi = m\mod n$ is a homomorphism with kernel $\langle n \rangle$. By the 
First Isomorphism Theorem, $\Z/\langle n\rangle \cong \Z_n$. 
\end{theorem}

If $\phi$ is a homomorphism from a finite group $G$ to $G'$, then $\|G\phi\|$ divides $\|G\}$ and $\|G'\|$. This follows from 
the First Isomorphism Theorem and the fact that $G\phi$ is a subgroup of $G'$ and Lagrange's Theorem. 

\begin{lemma}
The inner automorphism $\theta_x$ given by $g \goes x^{-1} g x$ is the identity automorphism if and only if $x\in C$ where $C$ is the center. 
\label{lemma1}\end{lemma}
If $x\in C$, $x^{-1} g x = x^{-1} x g = g$ and so $\theta_x = I$, the identity. 

If $\theta_x = I, x^{-1} g x = g \forall g \mbox{   and so  } x\in C$.

\begin{lemma}
The mapping $x \goes \theta_x$ of $G$ into the group of automorphisms of $G$ is a homomorphism
\end{lemma}
We already have that $\theta_{xy} = \theta_x\theta_y$ and so the result follows. 

By Lemma(\ref{lemma1}) the kernel of this homomorphism is the center $C$. The image is the subgroup of inner automorphisms. Therefore we obtain
\begin{theorem}
The group of inner automorphisms of a group $G \cong G/C$, where $C$ is the center of $G$.
\end{theorem}